\scanpage{7}

on est ramené à ceci.

\uline{Corollaire} Si $n\geq2$, $\underline{\mathcal L}_n$ est isomorphe par les hom \unsure{induits} à $\underline{\mathcal L}_2^{\otimes n'}$ si $n$ pair, à $\underline{\mathcal L}_2^{\otimes n'}\otimes\underline{\mathcal L}_3$ si $n$ impair (où $n'\geq1$ dans le premier cas, $n'\geq0$ dans le deuxième).

Mais on a plutôt intérêt de procéder directement ainsi : On est ramené aussitôt par platitude coh. au cas où $S=\operatorname{Spec}k$, on prend alors un élément $f$ de $L_n$ non dans $L_{n-1}$, un $g$ de $L_m$ non dans $L_{m-1}$, donc $fg$ est dans $\unsure{L_{n+m}|L_{n+m}^{0}}$ mais \unsure{provient} par \unsure{déf. de} $L_n\otimes L_m$.

\uline{Corollaire} $\mathcal E_3$ engendre $\Gamma_*(f,\underline L)$

Considérons les opérations de $\underline T$ sur les germes de fonctions rationnelles sur $C$ [reste à définir], on constate que \deleted{\unsure{$\underline T$}} ces opérations induisent des hom
\[
\underline{\mathcal L}\otimes\underline{\mathcal L}_n\longrightarrow\underline{\mathcal L}_{n+1}\qquad n\geq2
\]

\uline{Proposition 3} \inserted{Soit $s\in S$, $p$ l'exp. car. de $k(s)$.} L'homomorphisme $\underline{\mathcal L}\otimes\underline{\mathcal L}_n\unsure{\longrightarrow|-}\underline{\mathcal L}_{n+1}$ est un \uline{isomorphisme} (ou \deleted{\unsure{est}} un \unsure{épim.}) ssi $p$ est premier à $n$.

On est ramené au cas où $S=\operatorname{Spec}k$, alors c'est une vérification locale : si $D$ est une \deleted{\unsure{diff}} dérivation inv. non nulle en un \inserted{pt $x\in C$} et si $f\neq0$ sur $C$ a en $x$ un pôle d'ordre $n$,

\mspage{6}
